% TOP_PROBLEM: {"problemId": "problem.top500-omission-w041044-von-neumann-s-problem-for-ii-1-factors", "canonicalTitle": "Von Neumann's Problem for II_1 Factors", "releaseRank": 439, "primaryDomain": "analysis_pde", "primaryDomainLabel": "Analysis & PDE", "displayStatus": "open", "releaseStatus": "open"}
% !TeX program = lualatex
% Definition and sources only; this file is not an LLM solution attempt.
\documentclass[11pt]{article}
\usepackage[margin=25mm]{geometry}
\usepackage{fontspec}
\setmainfont{Latin Modern Roman}
\usepackage{amsmath,amssymb,mathtools}
\usepackage{unicode-math}
\setmathfont{Latin Modern Math}
\usepackage{xurl,hyperref}
\hypersetup{colorlinks=true,urlcolor=blue}
\setcounter{secnumdepth}{0}
\setlength{\parindent}{0pt}
\setlength{\parskip}{5pt}
\setlength{\emergencystretch}{3em}
\begin{document}
\section*{\texorpdfstring{Von Neumann's Problem for II\_1 Factors}{Von Neumann's Problem for II\_1 Factors}}\label{439}
\subsection{Definitions and mathematical statement}
A $\mathrm{II}_1$ factor is a von Neumann algebra with center ${\mathbb{C}}1$, a faithful normal tracial state, and no nonzero minimal projections. Here amenability is the usual injectivity/amenability property of von Neumann algebras. For a group $G$, $L(G)$ is the weak operator closure of the algebra generated by its left translation operators on $\ell^2(G)$.

Let $F_2$ be the free group on two generators. The question is whether every nonamenable $\mathrm{II}_1$ factor $M$ contains a unital von Neumann subalgebra isomorphic to $L(F_2)$, or whether a counterexample exists. Such a subalgebra is itself a factor, so this is also a subfactor-containment question. No finite-index hypothesis is added. The catalogue imposes no separability assumption, and none is silently supplied here.
\par\smallskip\textit{Formulation and concept references:} \href{https://math.vanderbilt.edu/peters10/problems.html}{[S1]}.

\subsection{Short English statement}
Does every nonamenable finite factor contain the von Neumann algebra of the free group on two generators?
\subsection{Sources}
\begin{itemize}
\item[S1]Open problems in operator algebras. Accessed 2026-09-01.
\url{https://math.vanderbilt.edu/peters10/problems.html}.
\textit{Formulation source.}
\end{itemize}
\end{document}
